%%%%%%%%%%%%%%%%%%%%%%%%%%%%%%%%%%%%%%%%%%%%
% 固定コマンド
%Don't touch the following setting
%
%%%%%%%%%%%%%%%%%%%%%%%%%%%%%%%%%%%%%%%%%%%%%



%%xcolor%%%%%%%%%%%%%%%%%%%%%%%%%%%%%%%%%%%%
\definecolor{darkgreen}{rgb}{0, 0.5, 0} %下の色指定で使う
\definecolor{whitesmoke}{rgb}{0.99, 0.99, 0.99} %下の色指定で使う
\newcommand{\fillme}{\xspace\textcolor{red}{XXX}\xspace}
\newcommand{\citeme}{\textcolor{red}{~\cite{FILL_ME_IN}}\xspace}

\newcommand{\red}{\color{red}}%red
\newcommand{\blue}{\color{blue}}
\newcommand{\green}{\color{teal}}
\newcommand{\purple}{\color{violet}}%violet
\newcommand{\black}{\color{black}}


\newcommand{\colorrow}{\rowcolor[rgb]{0.851,0.851,0.851}}
\newcommand{\colorcel}{\cellcolor[rgb]{0.851,0.851,0.851}}

%%Abbribiation%%%%%%%%%%%%%%%%%%%%%%%%%%%%%%%%%%%%%%%%%
\newcommand{\motivation}{\smallskip\noindent\textbf{Motivation. }}
\newcommand{\approach}{\smallskip\noindent\textbf{Approach. }}
\newcommand{\results}{\smallskip\noindent\textbf{Results. }}

\def\fig#1{Figure~\ref{#1}}
\def\tab#1{Table~\ref{#1}}
\def\eqn#1{Equation~\ref{#1}}
\def\sec#1{Section~\ref{#1}}
\def\Underline{\setbox0\hbox\bgroup\let\\\endUnderline}
\def\endUnderline{\vphantom{y}\egroup\smash{\underline{\box0}}\\}
\def\|{\verb|}

\newcommand{\ie}{\textit{i.e.,}\xspace}
\newcommand{\eg}{\textit{e.g.,}\xspace}

\newcommand{\etc}{\xspace\textit{etc.}\xspace}
\newcommand{\etal}{\xspace\textit{et al.}\xspace}
\newcommand{\aka}{\textit{a.k.a.,}\xspace}


%
\newcommand{\cnb}[3]{\fcolorbox{gray}{#1}{\bfseries\sffamily\scriptsize#2}{\sf\small\textit{\textcolor{red}{#3}}}}
\newcommand{\nb}[2]{\fcolorbox{gray}{yellow}{\bfseries\sffamily\scriptsize#1}{\sf\small\textit{\textcolor{red}{#2}}}}
\newcommand{\TODO}[1]{\cnb{yellow}{TODO:}{#1}}

\newcommand{\ASSIGN}[2]{\fcolorbox{gray}{yellow}{\bfseries\sffamily\scriptsize TODO (#1)}{\sf\small\textit{\textcolor{red}{#2}}}}

\newcommand{\doubt}[1]{\textcolor{orange}{#1}}
\newcommand{\reminder}[1]{\textcolor{blue}{#1}}




%Findings%%%%%%%%%%%%%%%%%%%%%%%%%%%%%%%%%%%%%%%%%%
% Define counters
\newcounter{findingctr}

% Define Finding command
\newcommand{\finding}[1]{%
  \noindent
  \refstepcounter{findingctr}%
  \textbf{Finding~\#\thefindingctr: #1}%
}

%Listing%%%%%%%%%%%%%%%%%%%%%%%%%%%%%%%%%%%%%%%%%%%%%%
\usepackage{listings}
\definecolor{backcolour}{rgb}{0.95,0.95,0.92}
\lstdefinelanguage{diff}{
  morecomment=**[f][\color{red}]{-},         % deleted lines 
  morecomment=**[f][\color{darkgreen}]{+},       % added lines
  moredelim=**[is][\bfseries]{@@}{@@},
}
\definecolor{backcolour}{rgb}{0.95,0.95,0.92}
\lstset{
  captionpos=t,
  language=java, % this is it !
  backgroundcolor=\color{backcolour},
  basicstyle={\ttfamily \scriptsize}, %書体を変える (scriptsizeで小さくしてます)
  %identifierstyle={\small}, %ちょっと字を小さくしてくれます
  %commentstyle
  % ={\ttfamily \scriptsize \color[rgb]{0,0.5,0}}, %コメントのフォント
  %keywordstyle={\small\bfseries}, %???
  %ndkeywordstyle={\small}, %???
  %stringstyle={\small\ttfamily}, %???
  frame=tb, %上下に線を引いてくれます
  breaklines=true, %行からはみ出た時に折り返してくれる
  columns=[l]{flexible}, %これがないと文字同士の間が妙に大きくなる．
  %ちなみにfullflexibleだとインデントを勝手に補完しちゃう分やりづらくなるので注意
  numbers=left, %左に行番号をつけたいときに使う
  xrightmargin=0em,
  xleftmargin=2em, %default:3zw.文字数分左余白を作る．もちろん負の値も可能．
  numberstyle={\scriptsize},
  stepnumber=1,
  numbersep=0.5em, %default:1zw zwは日本語用なので変更
  lineskip=-0.2ex, %行間
  tabsize=4,
  linewidth=\linewidth,
  breakindent = 6pt, %default:20pt (改行時のインデント)
  belowcaptionskip=0ex,
  aboveskip=-3pt,
  belowskip=-3pt,
  framesep=2pt, %フレームと中身（コード）の間隔
}
\lstdefinelanguage{commit}{ %コミットメッセージ用
  breakindent = 0pt,
  numbers=none,
  backgroundcolor=\color{white},
  frame=single,
  xleftmargin=3.5em,
  numbersep=0em,
  xrightmargin=1.5em,
}


%%tcolorbox%%%%%%%%%%%%%%
% \usepackage[skins]{tcolorbox}
% \newcommand{\conclusion}[1]{\begin{center}\begin{tcolorbox}[skin=widget, left=2mm,right=2mm,top=2mm,bottom=2mm,boxrule=0.3mm,arc=0mm,coltitle=black,colframe=black!99!white,colback=white!88!gray,width=(\linewidth),before=\hfill,after=\hfill]#1\end{tcolorbox}\end{center}}

\usepackage[many]{tcolorbox}  
\tcbuselibrary{listings,breakable}
\definecolor{main}{HTML}{D0D3D4}    % setting main color to be used
\definecolor{sub}{HTML}{D0D3D4}     % setting sub color to be used
\newtcolorbox{dbox}{
    left=2pt,right=2pt,top=2pt,bottom=2pt,
    enhanced, % for a fancier setting,
    boxrule = 0pt,% clearing the default rule
    % borderline = {0.2pt}{0pt}{main}, % outer line
    % borderline = {0.2pt}{1pt}{sub} % inner line
    enlarge top by=5pt,%     equivalent to mdframed 'skipabove'
    enlarge bottom by=3pt,%  equivalent to mdframed 'skipbelow'
  }
\tcbset{
    sharp corners,
    before skip = 0.1cm,    % add extra space before the box
    after skip = 0.01cm      % add extra space after the box
}


\def\summarybox#1#2{
\smallskip
\begin{tcolorbox}[
  enhanced,
  title=#1,
  colframe=darkgray,
  left=2mm,
  right=2mm,
]
    #2
\end{tcolorbox}
\medskip
}






% \setlength{\leftmargini}{1cm} % affects list environments
\renewenvironment{quote}
  {\list{}{\leftmargin=0.6cm \rightmargin=0.4cm}%
   \item\relax}
  {\endlist}





% Style 1: Classic monochrome code listing (IEEE/ACM style)
\lstdefinestyle{academicstyle}{
    basicstyle=\ttfamily\footnotesize,
    breaklines=true,
    frame=single,
    xleftmargin=2em,
    xrightmargin=2em,
    aboveskip=0.5em,
    belowskip=2em,
    abovecaptionskip=0em,      % キャプションの上の余白
    belowcaptionskip=0.5em,      % キャプションとコードの間の余白
    columns=flexible,
    keepspaces=true,
    language={},  % No language to prevent keyword highlighting
    % escapechar=@,
}




%Response letter%%%
\usepackage{tcolorbox}
\tcbuselibrary{skins,breakable}
\newtcolorbox{reviewcomment}[1][Comment]{  % デフォルトは"Conclusion"
    enhanced,
    frame hidden,
    borderline west={4pt}{0pt}{gray!50},
    colback=gray!10,
    boxrule=0pt,
    sharp corners,
    breakable,
    before upper={\parindent0pt\textbf{#1:}\par\vspace{0.3em}},
    fontupper=\itshape,
    left=10pt,
    right=10pt,
    top=5pt,
    bottom=5pt,
    before skip=1.5\baselineskip
}
\newcommand{\answer}[1]{\bigskip #1}

\newenvironment{src}{\color{gray}}{}
\newenvironment{rev}{\color{blue}}{}
\newcommand{\ADDED}[1]{\bigskip\noindent\textbf{\textsf{ADDED (#1):} }}
\newcommand{\DELETED}[1]{\bigskip\noindent\textbf{\textsf{DELETED (#1):} }}
\newcommand{\MODIFIED}[1]{\bigskip\noindent\textbf{\textsf{MODIFIED (#1):} }}

%%%%